\documentclass[12pt]{article}
\usepackage[T1]{fontenc}
\usepackage[utf8]{inputenc}
\usepackage{lmodern}
\usepackage{textcomp}
\usepackage{enumitem}
\usepackage{geometry}
\geometry{margin=1in}
\begin{document}
\begin{center}
Answer Key - Business Administration - Undergraduate Level Assessment 14 \
\small (Answers are ordered exactly as in the question paper.)
\end{center}

\section*{Multiple Choice}
\begin{enumerate}[label=\arabic*.]
\item \textbf{Answer:} B
\\ \textit{Options:} A) Egoism; B) Cognitive moral development; C) Power distance; D) Uncertainty avoidance
\\ \textit{Note:} The correct answer is "Cognitive moral development" because this theory, developed by Lawrence Kohlberg, outlines three distinct levels of moral reasoning-pre-conventional, conventional, and post-conventional-based on an individual's cognitive growth and capacity to evaluate ethical dilemmas.
\item \textbf{Answer:} C
\\ \textit{Options:} A) Safety; B) Belonging; C) Physiological; D) Esteem
\\ \textit{Note:} The physiological needs, such as food, water, and shelter, are considered the most basic in Maslow's hierarchy because they are fundamental for survival and must be satisfied before individuals can pursue higher-level psychological and self-fulfillment needs.
\item \textbf{Answer:} D
\\ \textit{Options:} A) Branding pyramid.; B) Brand value.; C) Brand names.; D) Brand personalities.
\\ \textit{Note:} The correct answer is "brand personalities" because it describes the practice of attributing human traits and characteristics to brands, allowing consumers to relate to them on a personal level.
\item \textbf{Answer:} C
\\ \textit{Options:} A) Replication.; B) Pre-code.; C) Reliability.; D) Convenience.
\\ \textit{Note:} The term "reliability" refers to the degree to which a scale produces consistent and stable ratings over time, making it the correct answer to describe the stability of measurement in this context.
\item \textbf{Answer:} D
\\ \textit{Options:} A) A group whose aim is transformational change; B) A group brought together to deliver training programmes; C) Team training for the purposes of advancing technology; D) Team building activities involving learning
\\ \textit{Note:} T-group training focuses on enhancing interpersonal skills and team dynamics through experiential learning activities, making it fundamentally about team building.
\end{enumerate}

\section*{True / False}
\begin{enumerate}[label=\arabic*.]
\setcounter{enumi}{5}
\item \textbf{Answer:} False
\\ \textit{Options:} A) True; B) False
\\ \textit{Note:} Price gouging typically refers to the unethical practice of raising prices excessively during emergencies or shortages, which is more commonly associated with B$_{2}$C contexts rather than B$_{2}$B, where pricing is generally governed by contracts and negotiations rather than opportunistic exploitation of market conditions.
\item \textbf{Answer:} False
\\ \textit{Options:} A) True; B) False
\\ \textit{Note:} The term DMU, or Decision-Making Unit, refers to the individuals or groups involved in the purchasing decision process rather than organizational characteristics like size and location.
\item \textbf{Answer:} True
\\ \textit{Options:} A) True; B) False
\\ \textit{Note:} The main purpose of licensing public relations practitioners is to ensure that professionals adhere to ethical standards and practices, ultimately safeguarding the interests and well-being of society by promoting responsible communication.
\item \textbf{Answer:} True
\\ \textit{Options:} A) True; B) False
\\ \textit{Note:} The Chase Jones model of issues management includes a comprehensive analysis step where the organization's strengths and weaknesses are assessed to identify how they can influence the management of emerging issues.
\item \textbf{Answer:} True
\\ \textit{Options:} A) True; B) False
\\ \textit{Note:} A scenario in scenario planning is considered an imagined sequence of future events because it represents a narrative that explores different potential outcomes and possibilities based on various assumptions and uncertainties in the business environment.
\end{enumerate}

\section*{Long Form Response}
\begin{enumerate}[label=\arabic*.]
\setcounter{enumi}{10}
\item \textbf{Answer:} An attitude transforms into an opinion when it is articulated and shared with others, often influenced by social interactions and the desire for validation. This process typically involves the individual reflecting on their internal beliefs and feelings, then framing these in a way that makes them communicable. In a business context, this transformation has significant implications; for instance, when employees express their opinions based on previously held attitudes, it can shape organizational culture, influence decision-making, and affect team dynamics. The ability to communicate opinions effectively is crucial, as it can lead to constructive dialogue or conflict, ultimately impacting overall business performance and innovation. Thus, understanding this transformation helps managers foster an environment where open communication leads to positive outcomes.
\\ \textit{Note:} correct because it clearly outlines the process of articulating and sharing internal attitudes, which transforms them into opinions, and emphasizes the importance of this transformation in influencing organizational dynamics and employee interactions in a business context.
\item \textbf{Answer:} Affirmative action in organizations refers to proactive measures aimed at supporting under-represented groups, such as women and minorities, in the workplace. These initiatives are designed to increase diversity, particularly in senior management roles, where historically, certain demographics have been disproportionately underrepresented. By implementing targeted recruitment, mentorship programs, and training opportunities, organizations can create a more inclusive environment that values diverse perspectives. The potential impact of these measures is significant, as a diverse leadership team can lead to enhanced decision-making, improved employee morale, and better overall organizational performance. Ultimately, affirmative action not only addresses historical inequities but also fosters a richer, more innovative workplace culture.
\\ \textit{Note:} correct because it accurately defines affirmative action as proactive initiatives aimed at enhancing diversity and inclusion, particularly for under-represented groups in senior management, through targeted strategies such as recruitment and mentorship.
\end{enumerate}

\vfill
\noindent\textit{End of Answer Key}
\end{document}